\sourcepage{1}
黎曼曲面起源于复分析，是处理多值函数问题的一种工具。这类多值函数的出现，是因为一个给定的全纯函数元素沿不同路径作解析延拓时，一般会得到该函数的不同分支。Riemann的想法是用复平面的一个多叶覆盖来代替函数的定义域。如果构造这一覆盖时，使平面上任一点上方的点数恰好等于该点处的函数元素个数，那么在这个“覆盖曲面”上，解析函数便成为单值函数。现在，忘掉这些曲面“铺展”在复平面（或黎曼球面）之上的事实，就得到抽象黎曼曲面的概念；它们可以看作单复变解析函数的自然定义域。

本章先讨论黎曼曲面的一般概念，随后从拓扑和解析两个角度研究覆盖空间。最后，把覆盖空间理论应用于解析延拓问题、代数函数黎曼曲面的构造、微分形式的积分以及线性微分方程的求解。
